\subsection{Product description}
\label{sec:product-desc}

In the tests, the integrated circuit is exposed to discharges while it is in normal operation.
The goal here is to identify functional failures and not hard-failures.
To operate correctly, the chip requires external components and cannot be tested in standalone.
The architecture of a typical application containing the chip, its external devices, and interaction with other modules is given in Fig. \ref{fig:system_architecture}.
In this application, the system is powered by the battery to the left.
The battery and the electronic module are connected with wires and cables.
Those cables constitute the most likely entry point for electrostatic discharges.

\begin{figure}[!h]
  \centering
  \includegraphics[width=\textwidth]{src/3/figures/architecture_system.pdf}
  \caption{Overview of the system architecture}
  \label{fig:system_architecture}
\end{figure}

Inside the electronic module (in blue), the \gls{asic} can be found (in red) along with decoupling capacitors for signals V\textsubscript{batt} and V\textsubscript{2p5}, and filtering devices.
The Common Mode Filter (CMF) and Pi filter are the two main devices being removed in the lighter configuration.

Inside the integrated circuit, two key parts of the circuit have been highlighted.
The primary supply is the function being studied in this chapter, and is powered by signal V\textsubscript{batt}.
Its detailed operation is provided later in this section, but overall it down converts the battery voltage.
On the output, it provides a regulated 2.5 V supply (signal V\textsubscript{2p5}).
V\textsubscript{2p5} is used further in the ASIC, to power the core functionnities of the device constituted of digital cells, low voltage analog functions, communication buses, etc.
The core of the circuit interacts with external equipments such as airbags for instance, and communicates with sensors or other modules.
Given this system architecture, it is obvious that the V\textsubscript{2p5} signal is critical for normal operation.
The disturbance of this signal could cause a propagation of issues way outside the integrated circuit, impacting other electronic modules and making some functionalities temporarily unavailable.

\subsubsection{Integrated supply function description}
\label{sec:supply-desc}

% Go down in the hierarchy to the primary supply and explain why testing this function
The next part of the study focuses on the primary supply function found inside the ASIC.
Preliminary analysis showed that this function is put at fault during ESD testing.
Appart from that, there are two motivations for testing this particular function.
The battery input V\textsubscript{batt} is an external pin, which is convenient to inject test stresses.
The regulated voltage output V\textsubscript{2p5} is also exposed externally, because the function requires a large capacitor for stabilisation.
This is very convenient to monitor externally the behavior of the primary supply.

% Main task
The main function of the regulation is to down-convert a battery voltage to a 2.5V regulated supply.
Its architecture is given in Fig. \ref{fig:monitored_function}.
Several blocks are involved for processing the battery supply.

\begin{figure}[!h]
  \centering
  \includegraphics[width=\textwidth]{src/3/figures/monitored_function.pdf}
  \caption{Architecture of the primary supply}
  \label{fig:monitored_function}
\end{figure}

% First block
The pre-regulator clamps the battery input voltage (V\textsubscript{batt}) that can reach up to 40 V and down to 9 V.
The clamped voltage is called V\textsubscript{clamp9}, has a current capability smaller than a few mA and is used for low-power functions.
The second output of the pre-regulator called V\textsubscript{pwr} provides a 12V clamped output with a larger current capability of tens of mA.

% Second block
The bandgap reference, connected downstream, is powered by V\textsubscript{clamp9}.
After a startup delay, it generates a 1.0 V voltage reference on V\textsubscript{ref1p0}.
This reference is stable accross a wide range of temperature, process variation and process mismatches.
The bandgap also outputs a 10 uA current reference on I\textsubscript{ref10u}, and the V\textsubscript{bgok} flag to signal it is ready for operation.

% Third major block
Finally, the \gls{ldo} regulator generates a stable 2.5V supply voltage on the external pin V\textsubscript{2p5}.
This output can deliver up to 20 mA.
V\textsubscript{2p5} is used further in the integrated circuit to power digital gates, like previously detailed.
A 100 nF decoupling capacitor is required on V\textsubscript{2p5} to absorb peak currents and help the regulator achieve stability.

% What are the minor blocks doing
The \textit{current mirrors} provide copied current values from the bandgap to the regulator, while offering much larger output impedance.
The \textit{enable} block mostly checks and waits for the bandgap to be properly started.
It then triggers a startup ramp-up sequence on the regulator.
